\section{A constructive proof from the search procedure}\label{ch09:sec:constructive}
Chapter~\ref{ch:matching} already contains almost everything we need for the sufficient direction. Recall the complete search of Section~\ref{ch08:sec:search} for a matching $M$, started at an unmatched student $x$. For each vertex it reaches, it records the predecessor, the vertex from which it was reached. It succeeds when it reaches an unused task, and fails when it has inspected every reached student without finding one. Chapter~\ref{ch:matching} proved a fact about each outcome.
\begin{itemize}
\item If the search succeeds, then following the predecessors back to $x$ gives an augmenting path (Section~\ref{ch08:sec:search}). Flipping $M$ along this path gives a matching with one more edge (Lemma~\ref{thm:augment}).
\item If the search fails, let $A$ be the set of students it reached, $x$ included, and let $B$ be the set of tasks it reached. Then $N(A)=B$ and $|B|=|A|-1$ (Proposition~\ref{ch08:prop:shortage}). Hence $|N(A)|=|A|-1<|A|$.
\end{itemize}
In the second case $A$ is a shortage set, so under \eqref{eq:hall} the second case can never happen. That is the whole idea of the proof.

\begin{proof}[Constructive proof that Hall's condition is sufficient]
Assume that the graph satisfies \eqref{eq:hall}. Start with the empty matching. As long as some student is unmatched, carry out the following round. Choose an unmatched student $x$ and run the complete search from $x$. If the search failed, the set $A$ of students it reached would satisfy $|N(A)|=|A|-1<|A|$ by Proposition~\ref{ch08:prop:shortage}, contradicting \eqref{eq:hall}. So the search reaches an unused task, and flipping the matching along the resulting augmenting path increases its number of edges by one.

Each edge of a matching contains exactly one student, and different edges contain different students, so a matching has at most $|L|$ edges. The matching starts with $0$ edges and gains one in every round, so there are at most $|L|$ rounds. The rounds stop only when no student is unmatched, so the final matching saturates $L$. (If $L=\varnothing$, there are no rounds, and the empty matching saturates $L$.)
\end{proof}

Section~\ref{sec:matching-trace} shows both outcomes. The graph with $N(a)=\{1,2\}$, $N(b)=\{1\}$ and $N(c)=\{2,3\}$ satisfies Hall's condition, and an unlucky start there, $a$--$1$ and $c$--$2$, in which $a$ takes the only task that $b$ accepts, is repaired along the augmenting path $b$--$1$--$a$--$2$--$c$--$3$, as the proof promises. When task~$3$ is removed from $c$'s options, the search from $b$ fails and reports the group $\{a,b,c\}$, which accepts only tasks $1$ and~$2$.

The procedure always ends with a certificate, a matching that saturates $L$ or a shortage set (Section~\ref{ch08:sec:shortage}), and Hall's theorem says that for each graph exactly one of the two exists: the necessary direction rules out a shortage set when a matching saturates $L$, and the sufficient direction, read as a contrapositive, says that when no matching saturates $L$, some group fails its test, and that group is a shortage set. This is Habit~7 (Section~\ref{ch08:sec:shortage}) carried to its end: every failure of the procedure produces a concrete obstruction, and a shortage set is the only kind of obstruction there is.
